\pagetitle{Portfolio and academic source record}{Materials, scholarship, and evidence boundaries}
\minorhead{Course readings supporting the one-page justification}
\begin{itemize}
\item Krajcik, J., \& Merritt, J. (2012). Engaging students in scientific practices: What does constructing and revising models look like in the science classroom? \emph{Science and Children}, 10--13.
\item Luehmann, A., Zhang, Y., Boyle, H., Tulbert, E., Merliss, G., \& Sullivan, K. (2024). Toward a justice-centered ambitious teaching framework: Shaping ambitious science teaching to be culturally sustaining and productive in a rural context. \emph{Journal of Research in Science Teaching, 61}, 319--357. \url{https://doi.org/10.1002/tea.21917}
\item Reiser, B. J., Novak, M., McGill, T. A. W., \& Penuel, W. R. (2021). Storyline units: An instructional model to support coherence from the students' perspective. \emph{Journal of Science Teacher Education, 32}(7), 805--829. \url{https://doi.org/10.1080/1046560X.2021.1884784}
\item Windschitl, M., \& Thompson, J. J. (2013). The modeling toolkit: Making student thinking visible with public representations. \emph{The Science Teacher}, 63--69.
\end{itemize}
\minorhead{Original teaching and implementation sources}
\begin{itemize}
\item Garcia, P., \& Aastha. (2026). \emph{Invisible Invaders: Monday, July 27, lesson and activity sequence}. Get Real! Science Camp. Original \href{https://docs.google.com/document/d/1jc_dSOKXVhSoSbvfbivXlgfYOxrrJ2L6KwbR7LQZCLc/edit}{Google Doc}; public \href{https://github.com/pzg8794/EDU486/blob/main/assignments/03-camp-unit-plan/monday-water-pathways-and-first-model.md}{EDU486 source}.
\item Garcia, P., \& Aastha. (2026). \emph{Invisible Invaders: Post-camp implementation revision addendum}. \href{https://github.com/pzg8794/EDU486/blob/main/assignments/03-camp-unit-plan/post-camp-implementation-revision-addendum.md}{EDU486 public-safe implementation record}. The co-facilitated camp, not this new design, generated the experience and feedback discussed here.
\item \emph{Water Test Investigation}. (2026). Collaborative EDE477 activity. \url{https://docs.google.com/document/d/1GpqZl8MVNCpH4dKFZKeCQcLOGir0535jPK4ayxFyxEs/edit}. Context for evidence-versus-claim distinctions, not the selected activity under revision.
\end{itemize}
\minorhead{What this project does and does not claim}
\begin{notebox}
\textbf{Original versus revised.} The full selected Activity 2 is reproduced in the appendix in black. Its original is co-authored by Piter and Aastha. Blue materials are proposed future lesson changes; they are not youth-produced products, implementation data, or independent evaluations of what worked.\par
\textbf{Scientific boundaries.} Paper proxies show one controllable transport mechanism. Field sampling shows collection and traceability. Suspected fragments are not chemically confirmed plastic, and neither model nor sample establishes exposure or harm. Student statements are paraphrased from public-safe records, not attributed to identified youth.\par
\textbf{Image and privacy boundaries.} The cropped field image is drawn from the existing July 2026 recap, excludes faces, and includes earlier AI-assisted layout work. It is not raw scientific evidence. There are no names, identifying student images, private placement records, or confidential data in this public artifact.\par
\textbf{AI process.} AI assisted with organizing the existing camp materials, annotated redesign, typography, and source comparisons under Piter's direction. Piter and Aastha retain credit for original camp work; Piter remains responsible for the redesign, review, and eventual submission.
\end{notebox}
\minorhead{Reuse and review}
This is designed as a transferable lesson module: preserve the unmodified original alongside the colored revision; retain the field-source links; adapt timing, materials, local site rules, and language routes for a new class. The assignment asks for a science activity from placement or freely available online: this is our authentic former camp lesson, available publicly at no cost. It is \emph{not} represented as a lesson from the current formal school placement or as instructor-approved under any narrower reading of that condition.
